\section{Light-Storage Protocols}
\label{sec:storage_protocols}

Light-storage protocols may be classified according to the physical mechanism used to control absorption and re-emission, the atomic level structure, or the parameter regime in which the interaction is operated. In the present discussion, a distinction is made between protocols based on engineered inhomogeneous broadening and optically controlled protocols~\cite{AWAVQMitCoSR_Rodriguez_}.

In engineered inhomogeneous broadening based protocols, the distribution of atomic resonance frequencies is prepared or modified so that the optical coherences accumulated by different frequency classes rephase at a later time. Controlled reversible inhomogeneous broadening (CRIB) and atomic-frequency-comb (AFC) storage belong to this category. In CRIB, the atomic detunings are reversed after absorption, causing the optical coherences to retrace their phase evolution and rephase. In AFC storage, the absorption line is prepared as a set of equally spaced peaks, and rephasing occurs after a delay fixed by their frequency spacing~\cite{Oqm_LvovskyEtAl_2009}. In their original two-level implementations, re-emission occurs after a fixed delay. For on-demand retrieval, the optical coherence must first be transferred to an additional long-lived state and subsequently restored by control pulses. These extended implementations therefore combine spectral engineering with optical control and may be regarded as hybrid protocols~\cite{AWAVQMitCoSR_Rodriguez_,Oqm_LvovskyEtAl_2009}.

In optically controlled memories, the transfer between the signal field and an atomic coherence is mediated by a strong classical field. The available protocols are then determined partly by the atomic level structure. In ladder-type systems, such as those used in Fast Ladder Memory (FLAME) and Off-Resonant Cascaded Absorption (ORCA), the signal and control fields couple a ground state to an intermediate state and subsequently to a higher excited state~\cite{ADToAEQMEBoaEQCD_RobertsonEtAl_2026, AWAVQMitCoSR_Rodriguez_}. In $\Lambda$-type systems, two long-lived lower states are coupled to a common excited state. This configuration supports electromagnetically induced transparency (EIT), off-resonant Raman storage, and Autler--Townes storage (ATS).

The operating regimes of these $\Lambda$-type protocols are distinguished mainly by the one-photon detuning, the control-field strength, and the timescale over which the control field is varied. Off-resonant Raman storage is performed at detunings much larger than the optical linewidth, such that the excited state is only weakly populated and an effective two-photon coupling is established between the lower states. This regime permits large storage bandwidths and has enabled broadband storage in alkali-metal vapours. Its implementation at the single-photon level is, however, limited by noise processes such as spontaneous Raman scattering and four-wave mixing when strong control fields are used~\cite{Qmeaara_HeshamiEtAl_2016,Oqm_LvovskyEtAl_2009}.

ATS is operated closer to resonance and relies on the splitting of the optical transition by a sufficiently strong control field. The signal pulse is absorbed into the resulting dressed-state resonances, and the corresponding atomic coherence is subsequently converted back into light by a second control pulse~\cite{AWAVQMitCoSR_Rodriguez_}. In contrast to adiabatic EIT storage, the interaction is governed by the coherent excitation and de-excitation of the split resonances rather than by propagation through a narrow transparency window.

Among the optically controlled protocols introduced above, EIT is of particular relevance because it is the protocol implemented in the experimental platform on which the present work is based. EIT operates in a near-resonant $\Lambda$-type system, where a strong control field induces destructive quantum interference and suppresses the absorption of the signal field. The resulting transparency window and steep dispersion enable the optical excitation to be mapped onto a collective ground-state coherence. Since this process determines both the physical origin and the spatial phase structure of the spin wave considered in the subsequent retrieval model, the EIT storage mechanism is examined in detail in the following subsection.

